\documentclass[12pt,a4paper]{article}

% ---------- Мова та шрифти ----------
\usepackage{fontspec}
\usepackage{polyglossia}
\setmainlanguage{ukrainian}
\setotherlanguage{english}

\setmainfont{DejaVu Serif}
\setsansfont{DejaVu Sans}
\setmonofont{DejaVu Sans Mono}

% ---------- Компонування ----------
\usepackage[a4paper,margin=2.2cm]{geometry}
\usepackage{graphicx}
\usepackage{float}
\usepackage{caption}
\usepackage{xcolor}
\usepackage{hyperref}
\usepackage{enumitem}
\usepackage{titlesec}
\usepackage{longtable}
\usepackage{array}
\usepackage{booktabs}
\usepackage{microtype}
\setlength{\emergencystretch}{3em}
\sloppy

\hypersetup{
    colorlinks=true,
    linkcolor=blue!50!black,
    urlcolor=blue!50!black,
    citecolor=blue!50!black,
    pdftitle={Звіт про десктоп-клієнт проєкту Remote File Folder Client},
    pdfauthor={Yana Utochkina}
}

\titleformat{\section}{\normalfont\Large\bfseries}{\thesection.}{0.6em}{}
\titleformat{\subsection}{\normalfont\large\bfseries}{\thesubsection}{0.6em}{}

\setlength{\parindent}{0pt}
\setlength{\parskip}{0.6em}

% Заглушка для скріншота: якщо файл #1 існує в img/screenshots/, вставляє
% реальне зображення; інакше -- підписану рамку з очікуваним шляхом до
% файлу, щоб документ завжди компілювався, навіть до додавання скріншотів.
\newcommand{\screenshot}[3]{%
    \begin{figure}[H]
        \centering
        \IfFileExists{img/screenshots/#1}{%
            \includegraphics[width=0.85\linewidth]{img/screenshots/#1}%
        }{%
            \fbox{\begin{minipage}[c][5.5cm][c]{0.85\linewidth}
                \centering
                \textit{Місце для скріншота}\\[6pt]
                \texttt{\small img/screenshots/#1}\\[6pt]
                \small додайте зображення за цим шляхом і перекомпілюйте
            \end{minipage}}%
        }
        \caption{#2}
        \label{#3}
    \end{figure}
}

\title{\textbf{Звіт про десктоп-клієнт\\[2pt] \large проєкту Remote File Folder Client\\(спрощений аналог Google Drive)}}
\author{Уточкіна Яна}
\date{Етап: десктоп-клієнт (macOS/SwiftUI) -- проєктування, використання та тестування \\ \today}

\begin{document}

\maketitle

\begin{abstract}
\noindent
Цей звіт документує \textbf{десктоп-клієнт macOS} проєкту \textbf{Remote File Folder Client} (\texttt{desktop/Sources/FileSpaceDesktop/}) -- спрощеного аналога Google Drive з веб- і десктоп-клієнтами, Go-бекендом, PostgreSQL та синхронізацією локальної теки. У звіті описано внутрішню структуру клієнта, розглянуто його прецеденти використання зі скріншотами та задокументовано покриття модульними тестами, включно з поведінкою сортування й фільтрації (сортування за \emph{``Edited by''}, фільтрація за \texttt{.cs}/\texttt{.jpg}).
\end{abstract}

\tableofcontents
\newpage

% =====================================================================
\section{Вступ}
% =====================================================================

Десктоп-клієнт -- один із двох фронтендів (другий -- веб-клієнт на React) до єдиного Go API та бази PostgreSQL, самостійно розгорнутих на домашньому Ubuntu-сервері за Cloudflare Tunnel. Усі прецеденти, які система підтримує через будь-який клієнт -- реєстрація, автентифікація, перегляд файлів із сортованими/фільтрованими атрибутами, перегляд вмісту \texttt{.java}/\texttt{.png}, вивантаження, завантаження, видалення -- доступні з десктоп-клієнта; крім того, він є \emph{єдиним} клієнтом, що реалізує безперервну фонову синхронізацію локальної теки, оскільки браузер не може постійно спостерігати за локальною текою.

% =====================================================================
\section{Архітектура десктоп-клієнта}
% =====================================================================

\subsection{Внутрішня структура}

Тека з вихідним кодом десктоп-клієнта, \texttt{desktop/Sources/\allowbreak FileSpaceDesktop/}, групує код так: SwiftUI-подання відокремлені від рушія синхронізації (\texttt{DirectoryWatcher}, \texttt{SyncViewModel} та \texttt{SyncState}/\texttt{FolderScanner} з \texttt{Manifest.swift}) і від спільного мережевого шару \texttt{APIClient}/\texttt{KeychainStore}:

\begin{itemize}[nosep]
    \item \textbf{Подання (Views)}: \texttt{LoginView}, \texttt{MainView} (таблиця файлів, сортування/фільтрація/видимість стовпців), \texttt{FilePreviewView} (\texttt{.java}/\texttt{.png}), \texttt{SyncPanelView} та \texttt{SyncConflicts\allowbreak SectionView} (керування синхронізацією теки та інтерфейс розв'язання конфліктів), \texttt{ContentView} (верхньорівнева навігація між станами ``не увійшов''/``увійшов'');
    \item \textbf{View-моделі}: \texttt{FilesViewModel} (список файлів, сортування/фільтрація, вивантаження/завантаження/видалення), \texttt{SyncViewModel} (цикл узгодження, виявлення/розв'язання конфліктів), \texttt{AppState} (стан сесії/автентифікації);
    \item \textbf{Мережа та сховище}: \texttt{APIClient} (REST-клієнт на \texttt{URLSession}, за протоколом \texttt{APIClientProtocol}, який використовується для тестування), \texttt{KeychainStore} (зберігання токенів через macOS Keychain), \texttt{Models.swift} (Codable DTO, що відповідають JSON-формату Go API);
    \item \textbf{Рушій синхронізації}: \texttt{DirectoryWatcher} (обгортка \texttt{FSEvents}), \texttt{Manifest.swift} (\texttt{FolderScanner}, \texttt{SyncState}, \texttt{SyncedFileEntry}).
\end{itemize}

Десктоп-клієнт звертається до Go API-сервера лише через HTTPS за публічним хостнеймом Cloudflare Tunnel -- тим самим, що використовує браузерний JavaScript веб-клієнта -- проходячи через серверні шари \texttt{Handler} $\to$ \texttt{Service} $\to$ \texttt{Repository} та, для вмісту файлів, інтерфейс \texttt{FileStorage}.

\subsection{Синхронізація}

Синхронізація -- найскладніша функція десктоп-клієнта і єдина, що не має аналога у веб-клієнті. За кожним тригером -- сповіщенням про зміну від \texttt{FSEvents} через \texttt{DirectoryWatcher} або періодичним опитуванням -- \texttt{SyncViewModel} викликає \texttt{POST /api/sync/diff}, надсилаючи \texttt{lastSyncedVersion} та маніфест (ім'я, розмір, sha256), обчислений \texttt{FolderScanner}; сервер повертає дії \texttt{download}/\texttt{delete\_local}/\texttt{conflict} разом із \texttt{newVersion}. Клієнт застосовує ці дії, а потім незалежно сканує свою локальну теку й вивантажує, повторно вивантажує або видаляє файли на основі власного збереженого \texttt{SyncState}.

\subsubsection*{Розв'язання конфліктів}

Дія \texttt{conflict} означає, що той самий файл змінився з обох боків від часу останньої синхронізації: сервер порівнює наданий клієнтом маніфест (ім'я, розмір, sha256) із власним записом і повертає дію \texttt{conflict} із розміром, часом модифікації та останнім редактором віддаленої сторони, коли хеші двох сторін справді не збігаються. На десктоп-клієнті \texttt{SyncViewModel} записує такі дії в опублікований список \texttt{conflicts} замість автоматичного розв'язання, а \texttt{SyncConflictsSectionView} (частина вже наявної панелі синхронізації) дозволяє користувачу обрати \textbf{Keep Local} (повторно вивантажити, перезаписавши віддалену копію), \textbf{Keep Remote} (перезаписати локальний файл) або \textbf{Keep Both} (зберегти локальний файл під перейменованою копією й прийняти віддалену версію).

% =====================================================================
\section{Прецеденти використання та скріншоти}
\label{sec:usecases}
% =====================================================================

У цьому розділі розглянуто прецеденти використання десктоп-клієнта, кожен зі скріншотом працюючого застосунку.

\subsection{Реєстрація та автентифікація}

\emph{Реєстрація} (рис.~\ref{fig:ss-register}) створює новий обліковий запис за логіном/паролем; \emph{автентифікація} (рис.~\ref{fig:ss-login}) потрібна для кожної іншої файлової операції, оскільки робоча область доступна лише після успішного входу, після чого \texttt{APIClient} зберігає отриману пару JWT access/refresh у macOS Keychain через \texttt{KeychainStore}.

\screenshot{01-register.png}{Екран реєстрації}{fig:ss-register}
\screenshot{02-login.png}{Екран входу}{fig:ss-login}

\subsection{Перегляд файлів: атрибути, стовпці, сортування, фільтрація}

Рис.~\ref{fig:ss-list} показує основну таблицю файлів з усіма видимими стовпцями атрибутів (ім'я, створено, змінено, автор вивантаження, автор редагування); стовпець імені приховати не можна, тоді як інші перемикаються через меню ``Columns'' (рис.~\ref{fig:ss-columns}).

\screenshot{03-file-list-overview.png}{Основний список файлів з усіма видимими стовпцями атрибутів}{fig:ss-list}
\screenshot{04-column-visibility.png}{Меню Columns: перемикання стовпця, відмінного від імені}{fig:ss-columns}

\textbf{Сортування:} список файлів сортується за \emph{``Edited by''} за зростанням або спаданням при кліку на заголовок цього стовпця, що циклічно перемикає \texttt{FilesViewModel.sortDirection} через none~$\to$~ascending~$\to$~descending~$\to$~none (рис.~\ref{fig:ss-sort}).

\screenshot{05-sort-edited-by.png}{Список файлів, відсортований за ``Edited by''}{fig:ss-sort}

\textbf{Фільтрація:} випадний список розширень, який формується з розширень, що реально присутні серед файлів користувача, фільтрує список лише до \texttt{.cs} або лише до \texttt{.jpg} (або назад до ``All'') через \texttt{FilesViewModel.extensionFilter} (рис.~\ref{fig:ss-filter}).

\screenshot{06-filter-extension.png}{Застосований фільтр за розширенням (\texttt{.cs} або \texttt{.jpg})}{fig:ss-filter}

\subsection{Перегляд вмісту файлу}

Клік на файл показує його вміст, якщо розширення -- \texttt{.java} (відображається як текст, рис.~\ref{fig:ss-java}) або \texttt{.png} (відображається як зображення, рис.~\ref{fig:ss-png}); файли з будь-яким іншим розширенням можна переглядати у списку й завантажувати, але не відображати їхній вміст.

\screenshot{07-preview-java.png}{Файл \texttt{.java}, показаний як текст}{fig:ss-java}
\screenshot{08-preview-png.png}{Файл \texttt{.png}, показаний як зображення}{fig:ss-png}

\subsection{Вивантаження, завантаження, видалення}

Рис.~\ref{fig:ss-upload} показує вивантаження в процесі -- великі файли розбиваються на chunk'и, щоб жоден окремий HTTP-запит не перевищував ліміт розміру тіла безкоштовного тарифу Cloudflare Tunnel, а \texttt{uploadProgress} відображає кількість надісланих байтів по кожному chunk'у. Рис.~\ref{fig:ss-download} та \ref{fig:ss-delete} показують завантаження файлу в теку Downloads користувача та видалення файлу з віддаленої робочої області.

\screenshot{09-upload.png}{Вивантаження файлу в процесі}{fig:ss-upload}
\screenshot{10-download.png}{Завантаження файлу}{fig:ss-download}
\screenshot{11-delete.png}{Видалення файлу}{fig:ss-delete}

\subsection{Синхронізація локальної теки}

Рис.~\ref{fig:ss-sync} показує панель синхронізації з обраною текою й активною синхронізацією: \texttt{DirectoryWatcher} генерує сповіщення про зміни на основі \texttt{FSEvents}, а \texttt{SyncViewModel} також періодично опитує сервер, виконуючи цикл узгодження, описаний у \S2.2. Рис.~\ref{fig:ss-conflict} показує ту саму панель, коли сервер повідомляє про справжній конфлікт: файл показано разом із його локальним/віддаленим розміром та редактором, а також діями Keep Local / Keep Remote / Keep Both.

\screenshot{12-sync-panel.png}{Панель синхронізації теки, активна синхронізація}{fig:ss-sync}
\screenshot{13-sync-conflict.png}{Панель синхронізації теки з виявленим конфліктом та варіантами його розв'язання}{fig:ss-conflict}

% =====================================================================
\section{Модульне тестування}
\label{sec:testing}
% =====================================================================

\subsection{Налаштування тестів}

Тестова ціль десктоп-клієнта, \texttt{FileSpaceDesktopTests}, використовує \texttt{XCTest} (запускається через \texttt{swift test} з \texttt{desktop/} або через \texttt{make desktop-test}), а не новіший модуль \texttt{swift-testing}, оскільки в цьому середовищі встановлено лише Command Line Tools від Xcode, які не постачають модуль \texttt{Testing}. Мережеві або залежні від ОС співучасники замінюються in-memory заглушками, тож набір тестів виконується офлайн і детерміновано:

\begin{itemize}[nosep]
    \item \texttt{SyncViewModel} і \texttt{FilesViewModel} залежать від \texttt{APIClientProtocol}, а не від конкретного \texttt{APIClient}, тож тести підставляють безмережевий \texttt{FakeAPIClient} замість \texttt{URLSession};
    \item тести \texttt{FolderScanner} і стану синхронізації використовують реальні тимчасові каталоги на диску, а не заглушки файлової системи, оскільки код, що тестується, \emph{сам є} взаємодією з файловою системою;
    \item \texttt{KeychainStoreTests} працюють із реальним macOS Keychain (in-memory заміни для нього немає), використовуючи виділені ключі, які очищаються в \texttt{tearDown}.
\end{itemize}

\subsection{Покриття за областями}

Таблиця~\ref{tab:tests} підсумовує набір модульних тестів.

\small
\begin{longtable}{@{}p{4.8cm}p{1.2cm}p{9.2cm}@{}}
\caption{Набори модульних тестів у \texttt{desktop/Tests/\allowbreak FileSpaceDesktopTests/}}
\label{tab:tests} \\
\toprule
\textbf{Клас тестів} & \textbf{К-сть} & \textbf{Що покриває} \\
\midrule
\endhead
\texttt{FileRecord\-Decoding\-Tests} & 2 & Декодування \texttt{FileRecord} з точного JSON-формату бекенду, включно з перейменуванням ключа \texttt{extension} $\to$ \texttt{extensionName} \\
\texttt{SyncStateTests} & 1 & Round-trip \texttt{SyncState}/\texttt{SyncedFileEntry} через кодування у JSON \\
\texttt{FolderScannerTests} & 4 & Сканування локальної теки: обчислення хешу, пропуск прихованих файлів/підкаталогів та коректне розрізнення ``неможливо прочитати'' від ``не існує'' \\
\texttt{PreviewSupportTests} & 1 & Лише \texttt{.java}/\texttt{.png} вважаються придатними для перегляду \\
\texttt{KeychainStoreTests} & 5 & Встановлення/читання/перезапис/видалення токенів у реальному macOS Keychain з ізоляцією ключів \\
\texttt{APIClientTests} & 6 & Декодування запитів/відповідей, порядок і прогрес chunked-вивантаження, вибір між одноразовим і chunked-вивантаженням, поведінка повтору запиту після 401 з оновленням токена \\
\texttt{SyncViewModelTests} & 11 & Цикл узгодження: обробка нових/змінених/локально видалених файлів, надсилання маніфесту, виявлення/фіксація конфліктів та всі три способи розв'язання (Keep Local/Remote/Both) \\
\texttt{FilesViewModelTests} & 7 & Сортування за \emph{``Edited by''} (за зростанням, за спаданням, без урахування регістру) та фільтрація за розширенням (\texttt{.cs}, \texttt{.jpg}, без фільтра), включно з одночасним застосуванням сортування й фільтра \\
\bottomrule
\textbf{Разом} & \textbf{37} & \\
\end{longtable}

\subsection{Покриття сортування та фільтрації}
\label{sec:variant}

Список файлів сортується за \textbf{``Edited by''} (за зростанням і за спаданням) та фільтрується до \textbf{всіх файлів}, лише файлів \texttt{.cs} або лише файлів \texttt{.jpg}. Ця логіка міститься в \texttt{FilesViewModel.displayedFiles}, \texttt{toggleSort(\_:)} та \texttt{extensionFilter} і перевіряється тестами з \texttt{FilesViewModelTests.swift}:

\begin{itemize}[nosep]
    \item[] \texttt{testSortByEditedByAscending}
    \item[] \texttt{testSortByEditedByDescending}
    \item[] \texttt{testSortByEditedByIsCaseInsensitive}
    \item[] \texttt{testNoFilterShowsAllExtensions}
    \item[] \texttt{testFilterByCsExtensionShowsOnlyCsFiles}
    \item[] \texttt{testFilterByJpgExtensionShowsOnlyJpgFiles}
    \item[] \texttt{testFilterAndSortByEditedByComposeTogether}
\end{itemize}

\texttt{FilesViewModel} залежить від \texttt{APIClientProtocol}, а не від конкретного, прив'язаного до мережі \texttt{APIClient} (той самий принцип, що вже використовує \texttt{SyncViewModel}), і саме це дозволяє цим тестам підставляти заглушку та виконуватися офлайн.

\begin{verbatim}
Test Suite 'All tests' passed at 2026-09-30 10:53:05.
  Executed 37 tests, with 0 failures (0 unexpected)
  in 0.190 (0.197) seconds
\end{verbatim}

% =====================================================================
\section{Висновки}
% =====================================================================

Десктоп-клієнт реалізує кожен прецедент, потрібний проєкту, який нативний застосунок macOS може змістовно додати порівняно з веб-клієнтом: повний перегляд файлів із сортуванням за ``Edited by'' та фільтрацією за \texttt{.cs}/\texttt{.jpg}, перегляд вмісту \texttt{.java}/\texttt{.png}, вивантаження/завантаження/видалення, а також -- унікально серед обох клієнтів -- безперервну фонову синхронізацію теки, включно з виявленням конфліктів на основі маніфесту та трибічним (Keep Local/Remote/Both) інтерфейсом розв'язання. З погляду тестування, основні співучасники клієнта (мережа, сканування теки, стан синхронізації, доступ до Keychain, цикл узгодження та сортування/фільтрація списку файлів) мають модульне покриття. Усі 37 тестів набору проходять успішно.

\end{document}
